% Figure: Brayton net specific work against pressure ratio for three values of
% the temperature ratio tau = T3/T1. Each curve peaks at an intermediate
% pressure ratio, the specific-work optimum r_p = tau^(gamma/(2(gamma-1))).
% Normalised: w_net / (c_p T_1) = tau(1 - x^-1) - (x - 1), with x = r_p^0.286.
% The dimensionless group x sweeps; we plot against r_p directly.
\begin{figure}[htbp]
\centering
\begin{tikzpicture}
\begin{axis}[
  width=12cm, height=8cm,
  xlabel={pressure ratio $r_p$},
  ylabel={normalised net work $w_{\text{net}}/(c_p T_1)$},
  xmin=1, xmax=45, ymin=0, ymax=2.6,
  axis lines=left,
  legend pos=north east, legend cell align=left,
  tick label style={font=\scriptsize}, label style={font=\small},
  legend style={font=\scriptsize}, clip=false,
]
% tau = 4: w/(cp T1) = 4*(1 - x^-1) - (x-1), x = r_p^0.285714
\addplot[engblue, very thick, samples=120, domain=1:45]
  {4*(1 - (x^0.285714)^(-1)) - (x^0.285714 - 1)};
\addlegendentry{$\tau = T_3/T_1 = 4$}
% tau = 5
\addplot[engred, very thick, samples=120, domain=1:45]
  {5*(1 - (x^0.285714)^(-1)) - (x^0.285714 - 1)};
\addlegendentry{$\tau = 5$}
% tau = 6
\addplot[engamber, very thick, samples=120, domain=1:45]
  {6*(1 - (x^0.285714)^(-1)) - (x^0.285714 - 1)};
\addlegendentry{$\tau = 6$}
% peak markers: r_p = tau^(gamma/(2(gamma-1))) = tau^1.75
% tau=4: 4^1.75 = 11.3 ; tau=5: 5^1.75=16.7 ; tau=6: 6^1.75=22.9
\addplot[engblack, only marks, mark=*, mark size=1.6pt] coordinates {(11.3,1.07)};
\addplot[engblack, only marks, mark=*, mark size=1.6pt] coordinates {(16.7,1.53)};
\addplot[engblack, only marks, mark=*, mark size=1.6pt] coordinates {(22.9,2.02)};
\node[font=\scriptsize, anchor=south] at (axis cs:11.3,1.10) {$r_p^\star$};
\end{axis}
\end{tikzpicture}
\caption[Brayton specific work versus pressure ratio]{Normalised net specific
work of the air-standard Brayton cycle against pressure ratio, for three values
of the temperature ratio $\tau = T_3/T_1$. Each curve rises from zero at
$r_p = 1$, where the cycle does no work, climbs to a peak, and falls back toward
zero at the ceiling ratio $r_p = \tau^{\gamma/(\gamma-1)}$, where the compressor
exit reaches the turbine inlet temperature and no heat can be added. The peak
(marked $r_p^\star$) sits at $r_p^\star = \tau^{\gamma/(2(\gamma-1))}$, the
geometric mean of unity and the ceiling ratio, and it moves to higher pressure
ratio as the turbine inlet temperature rises. The specific-work optimum lies
well below the pressure ratio that maximises thermal efficiency, which is the
split a designer resolves by weighing engine size against fuel cost.}
\label{fig:vol04ch03:brayton-specific-work}
\end{figure}
